Vanishing of one Gauduchon degree forces a positive current to vanish.

\begin{lemma}\label{lem:zero-gauduchon-degree}
\index{Gauduchon degree!vanishing}
Let $\theta$ be a smooth real $(1,1)$-form with $\PSH(X,\theta)\ne\varnothing$.
If
\[
 \int_X\theta\wedge\gamma^{n-1}=0
\]
for one Gauduchon metric $\gamma$, then $\theta=\ddc f$ for a smooth real
function $f$. Conversely, a smoothly $\ddc$-exact form has zero degree
against every Gauduchon metric. The conclusion also holds for $n=1$, with $\gamma^{n-1}=1$.
\end{lemma}
\Needspace{8\baselineskip}
\begin{proof}
Choose $u\in\PSH(X,\theta)$ and put $S=\theta+\ddc u\geq0$. Stokes' theorem gives $\int_XS\wedge\gamma^{n-1}=0$. Since the trace measure controls all coefficient measures of $S$, one has $S=0$. Taking the trace of $\ddc u=-\theta$ with respect to $\omega$ gives an elliptic equation with smooth right-hand side. Elliptic regularity makes $u$ smooth, so $f=-u$ works. The converse is Stokes' theorem. For $n=1$, the trace measure is $S$ itself and the same argument applies.
\end{proof}

For a smooth real $(1,1)$-form $\theta$, put $a=\kappa_\omega(\theta)$. By \cref{cor:gauduchon-threshold},
$\theta-a\omega$ admits a positive current. Hence
\cref{lem:zero-gauduchon-degree} shows that the infimum defining
$\kappa_\omega(\theta)$ is attained by a smooth Gauduchon metric if and
only if
\[
 \theta-a\omega=\ddc f\quad\text{for some }f\in C^\infty(X,\mathbb R).
\]
In that case every Gauduchon metric attains it. The positive current at
the threshold always exists; a smooth minimizing test is much more restrictive.

\section{Compactness and duality of Gauduchon currents}
\label{sec:gauduchon-linear-duality}

The weak closure of normalized smooth Gauduchon forms supplies a minimizing test even when no smooth form attains the threshold. The normalization gives compactness by bounding the coefficient measures. Let $\mathscr G_\omega$ be the set of positive $(n-1,n-1)$-currents $T$ on $X$ satisfying
\index{Gauduchon current!normalized tests}
\[
 \ddc T=0,\qquad\int_X\omega\wedge T=1.
\]
These are the normalized Gauduchon currents. Let $\mathscr C_\omega$
be the weak closure of the smooth strictly positive members of
\index{Gauduchon current!closure of smooth tests}
$\mathscr G_\omega$. The root calculation for \cref{thm:lamari} identifies those members with normalized powers of
Gauduchon metrics.

\begin{proposition}\label{prop:gauduchon-linear-compact-dual}
The sets $\mathscr G_\omega$ and $\mathscr C_\omega$ are nonempty,
compact and convex. For every smooth real $(1,1)$-form $\theta$,
\begin{equation}\label{eq:gauduchon-linear-minimum}
 \kappa_\omega(\theta)
   =\min_{T\in\mathscr C_\omega}\int_X\theta\wedge T.
\end{equation}
Moreover, for $T\in\mathscr G_\omega$, membership in
$\mathscr C_\omega$ is equivalent to
\begin{equation}\label{eq:gauduchon-linear-polar}
 \int_X\theta\wedge T\geq0
 \quad\text{for every smooth real }\theta
       \text{ with }\PSH(X,\theta)\ne\varnothing.
\end{equation}
\end{proposition}
\begin{proof}
Gauduchon's theorem supplies a smooth positive member after normalization.
The fixed trace mass bounds the total variation of all coefficient
measures in finitely many local frames. Weak compactness of finite
measures therefore gives compactness of the family of positive currents
with this trace. Positivity, the trace normalization and $\ddc T=0$
pass to limits. On this bounded set, testing against smooth forms or
against continuous coefficient forms gives the same topology, by
uniform approximation. It is metrizable since the continuous test
space is separable. Thus the compactness can also be used sequentially.
The defining conditions of $\mathscr G_\omega$ are convex. Its smooth
strictly positive members form a convex set, so their closure
$\mathscr C_\omega$ is compact and convex as well.

Pairing with a fixed smooth $\theta$ is continuous. Its minimum on the
closure of the smooth tests is therefore their infimum, proving
\eqref{eq:gauduchon-linear-minimum}. The Lamari criterion shows that
\eqref{eq:gauduchon-linear-polar} holds for every smooth positive test,
and hence for every member of its weak closure.

Conversely, if $T\notin\mathscr C_\omega$, separate this point from the
compact convex set $\mathscr C_\omega$ in the weak space of currents.
A continuous linear functional for that weak topology is evaluation
against a smooth form. After changing its sign if necessary, there is
a smooth real $(1,1)$-form $\beta$ such that
\[
 \int_X\beta\wedge T
 <\min_{R\in\mathscr C_\omega}\int_X\beta\wedge R
 =\kappa_\omega(\beta).
\]
By \cref{cor:gauduchon-threshold}, the smooth form
$\beta-\kappa_\omega(\beta)\omega$ admits a positive current.
Its pairing with $T$ is negative because $\int_X\omega\wedge T=1$.
This contradicts \eqref{eq:gauduchon-linear-polar} and proves the converse.
\end{proof}

\begin{corollary}\label{cor:gauduchon-threshold-continuity}
The function $\kappa_\omega$ is concave, positively homogeneous and
invariant under addition of $\ddc f$ for $f\in C^\infty(X,\mathbb R)$. Let $\theta$ and $\eta$ be smooth real $(1,1)$-forms and let $\varepsilon\geq0$. If
$-\varepsilon\omega\leq\theta-\eta\leq\varepsilon\omega$, then
\[
 |\kappa_\omega(\theta)-\kappa_\omega(\eta)|\leq\varepsilon.
\]
If smooth real $(1,1)$-forms $\theta_j$ converge uniformly to a smooth form $\theta$
and $T_j\in\mathscr C_\omega$ minimize their pairings, every weak
cluster point of $(T_j)$ minimizes the pairing with $\theta$.
\end{corollary}
\begin{proof}
Formula \eqref{eq:gauduchon-linear-minimum} is an infimum of linear
functionals, which proves concavity and positive homogeneity.
Gauge invariance follows from $\ddc T=0$. The normalized trace gives
\[
 \left|\int_X(\theta-\eta)\wedge T\right|\leq\varepsilon
 \quad(T\in\mathscr C_\omega),
\]
and taking infima proves the estimate. For the last assertion, pass to
a subsequence $T_j\rightharpoonup T\in\mathscr C_\omega$.
The estimate and uniform convergence of the backgrounds give
\[
 \int_X\theta\wedge T
 =\lim_j\int_X\theta_j\wedge T_j
 =\lim_j\kappa_\omega(\theta_j)=\kappa_\omega(\theta).
\]
Thus $T$ is a minimizer.
\end{proof}

The closure $\mathscr C_\omega$ cannot in general be replaced by all of
$\mathscr G_\omega$. The exceptional-curve example in
\cref{ex:gauduchon-current-smoothing-obstruction} separates the two
sets by a smooth closed form. Testing against all of $\mathscr G_\omega$
will instead detect smooth positivity and nefness in
\cref{thm:smooth-cone-current-duality}.
